\scanpage{408}
\mspage{322}
% archive sheet number: 407

% 1–2
qu'on peut trouver $i\xrightarrow{\unsure{\rho|s}}i'$ dans $I$, tel que la flèche composée $\varphi(i_{\alpha_0})\xrightarrow{w}\varphi(i)\xrightarrow{\varphi(\rho)}\varphi(i')$
% 3–5
soit de la forme $\varphi(\tau)$. Donc quitte à remplacer $i$ par $i'$, et $v_{\alpha_0}$ par \unsure{la|le} \unsure{composé|composée} \inserted{$v'_{\alpha_0}:$} $\unsure{j|J}_{\deleted{\unsure{\beta}}\alpha_0}\xrightarrow{v_{\alpha_0}}\varphi(i)\xrightarrow{\varphi(\rho)}\varphi(i')$, \deleted{\unsure{OPS|SPS}}
% 6–10
\deleted{que} \unsure{OPS|SPS} que $v_{\alpha_0}u_{\alpha_0}$ est de la forme $\varphi(\tau)$. Cela nous donne donc en principe \unsure{l'}objet $i=i_\alpha$, et les flèches $v_{\alpha_0}:$ \inserted{$j_{\alpha_0}\to\varphi(i_{\unsure{\alpha|\tau}})$} et $\tau_{\alpha,\alpha_0}:i_{\alpha_0}\unsure{-|\to}i_\alpha$, de façon à satisfaire à la première compati\deleted{\unsure{ti}}bilité \unsure{dans|de} \unsure{(*)|(a)},

% 11–15
Il reste à construire $j_\alpha,\tau'_{\alpha,\alpha_0},u_\alpha$ de façon à satisfaire à la deuxième compati\deleted{\unsure{o}}bilité $(*)$. Or $J'$ étant cofinal \unsure{dans|sur} $J$, il existe $j_\alpha$ dans $J'$ et une flèche $u_\alpha:\varphi(i_\alpha)\longrightarrow\deleted{\unsure{J}}j_\alpha$.

% 16–20
D'autre part, comme $J'\unsure{-|\to}J$ est pl. fidèle, la flèche composée $j_{\alpha_0}\xrightarrow{v_{\alpha_0}}\varphi(i_\alpha)\xrightarrow{u_\alpha}j_\alpha$ dans $\unsure{\mathbb{J}|J|I}$, entre objets de $J'$, provient d'une unique flèche $\tau'_{\unsure{\alpha_0|\alpha},\alpha_0}:j_{\alpha_0}\longrightarrow j_\alpha$ dans $J'$.
